\section{程序说明与核心代码}

实现语言为 Python。参数由 \texttt{parameters.csv} 冻结，求解器不另写一套常数。在线程序只使用指令反馈，不导入含源真值的校验模块。回归测试 $168$ 项通过。问题3入口默认 \texttt{GREEDY\_FAST}，问题4入口默认 HEX37 与 \texttt{PREFIX\_A}；正式测试入口在写作时未启用。

主要文件：\texttt{params.py} 读入半径、误差与速度；\texttt{geometry.py} 给出楔形半平面、顶点枚举、最远对与最小包围圆；\texttt{localization.py} 给出理论扇半径、$J(q)$ 与第二点分层采样；\texttt{q3\_state.py}、\texttt{q3\_runner.py}、\texttt{q3\_greedy\_policy.py}、\texttt{q3\_fast\_mode.py}、\texttt{q3\_safe.py} 构成问题3闭环；\texttt{q4\_cover.py}、\texttt{q4\_localize.py}、\texttt{q4\_policy.py}、\texttt{q4\_runner.py} 构成问题4闭环。完整源文件置于支撑材料，此处仅列出与正文命题直接对应的片段。

\subsection{理论扇半径}

\begin{lstlisting}
def theoretical_sector_rho(r_max=R_MAX, eps_deg=EPS_DEG) -> float:
    return float(r_max / (2.0 * np.cos(np.deg2rad(eps_deg))))
\end{lstlisting}

\subsection{HEX37 覆盖半径与冻结路径}

\begin{lstlisting}
HEX_STEP = 800.0
HEX_RING = 3
HEX_DELTA = HEX_STEP / np.sqrt(3.0)
HEX_E1 = np.array([HEX_STEP, 0.0])
HEX_E2 = np.array([HEX_STEP / 2.0, HEX_STEP * np.sqrt(3.0) / 2.0])
# 36 edges * 800 m = 28800 m. Offline Hamilton path from the origin.
HEX37_ROUTE_QR = (
    (0, 0), (-1, 0), (-1, 1), (-2, 1), (-3, 1), (-3, 0), (-2, 0),
    (-2, -1), (-1, -2), (-1, -1), (0, -1), (0, -2), (0, -3), (1, -3),
    (1, -2), (2, -3), (3, -3), (2, -2), (1, -1), (1, 0), (0, 1),
    (-1, 2), (-2, 2), (-3, 2), (-3, 3), (-2, 3), (-1, 3), (0, 3),
    (0, 2), (1, 2), (1, 1), (2, 1), (3, 0), (2, 0), (2, -1), (3, -2),
    (3, -1),
)
\end{lstlisting}

\subsection{光学方格边长}

边长取 $1.4\times 20=28\,\mathrm{m}$，半对角线 $14\sqrt{2}<20$。定位域按主轴旋转后与方格求交，每个非空交集的包围中心作为一个清除点。该片段覆盖连续多边形，不抽样顶点代替区域。
